%\documentclass[a4paper, 14pt]{article}
\documentclass[a4paper, 14pt]{extarticle} 

\usepackage[utf8]{inputenc}
\usepackage[T2A]{fontenc}
\usepackage[english,russian]{babel}
% \fontsize{14}{17}

% \usepackage[none]{hyphenat}  % Полное отключение переносов

\usepackage{setspace}
%\полуторный интервал
% \onehalfspacing
\usepackage[left=30mm, top=20mm, right=15mm, bottom=20mm, nohead, nofoot]{geometry}
\usepackage{amsmath,amsfonts,amssymb, gensymb, tikz, graphicx}
\usepackage{fancybox,fancyhdr}
\usepackage{multicol}
\usepackage{indentfirst}
\usepackage{tabto}
\usepackage{titlesec}
\usepackage{listings}
\usepackage{fancyhdr}
\usepackage{booktabs}
\usepackage{float}
\usepackage{tocloft} % для настройки оглавления
\usepackage{titlesec} % для настройки заголовков
\usepackage{hyperref} % переносит длинные ссылки
\pagestyle{fancy}
\fancyhf{}  % Очистить стандартные колонтитулы
\renewcommand{\headrulewidth}{0pt} % Убираем линию под верхним колонтитулом
\renewcommand{\footrulewidth}{0pt} % Убираем линию над нижним колонтитулом (если нужно)
\cfoot{\thepage}  % Номер справа внизу
% \pagestyle{fancy}
% \fancyhf{}
% \fancyhead[R]{A. Khalitov}
% \fancyfoot[R]{\thepage}
% \fancyhead[L]{Срез знаний за 7 класс}
\usepackage[ruled,vlined]{algorithm2e}
\usepackage{algorithm, algorithmic, algpseudocode}
\setcounter{page}{4}
% \headsep=10mm
\usepackage{xcolor}
% \usepackage[style=gost-numeric]{biblatex} % для оформленния ссылок на литературу по ГОСТу
\usepackage{hyperref}
\hypersetup{colorlinks=true, allcolors=black}
\newcommand{\lr}[1]{\left( {#1} \right)}

% Включаем отступ после \subsection
\makeatletter
\patchcmd{\@startsection}{\@afterindentfalse}{\@afterindenttrue}{}{}
\makeatother

\usepackage{csquotes} % требуется biblatex
\usepackage[style=numeric,sorting=none]{biblatex}
\addbibresource{sources.bib}

% Настройка — оставить только название источника в списке литературы
\AtEveryBibitem{%
	\clearfield{author}%
	\clearfield{editor}%
	\clearfield{publisher}%
	\clearfield{year}%
	\clearfield{journal}%
	\clearfield{note}%
}

% Переопределим шаблоны, чтобы показывать только заголовок
\DeclareFieldFormat{title}{#1}
\DeclareBibliographyDriver{book}{\printfield{title}}
\DeclareBibliographyDriver{article}{\printfield{title}}
\DeclareBibliographyDriver{online}{\printfield{title}}

% \usepackage[backend=biber, style=numeric, sorting=none]{biblatex}

% % Настраиваем отображение записей в библиографии
% \DeclareFieldFormat[article,book]{title}{#1} % Название выводится первым
% \renewbibmacro{in:}{} % Убираем "In:" для статей

% % Настраиваем отображение ссылок
% \DeclareFieldFormat{labelnumberwidth}{[#1]} % Формат ссылок: [1], [2], ...

% Отступы
% Абзацный отступ
\setlength{\parindent}{1.25cm}
% Настройка отступов для section
\titleformat{\section}
{\normalfont\Large\bfseries}{\thesection}{1em}{} % {формат}{метка}{отступ от метки}{перед заголовком}[после заголовка]
\titlespacing*{\section}{1.25cm}{3.5ex}{2.3ex} % {левый отступ}{отступ сверху}{отступ снизу}

% Настройка отступов для subsection
\titleformat{\subsection}
{\normalfont\large\bfseries}{\thesubsection}{1em}{}
\titlespacing*{\subsection}{1.25cm}{2.5ex}{1.5ex}

% Настройка отступов для subsubsection
\titleformat{\subsubsection}
{\normalfont\large\bfseries}{\thesubsubsection}{1em}{}
\titlespacing*{\subsubsection}{1.25cm}{2.5ex}{1.5ex}
% \titleformat{\section}{\normalfont\normalsize\bfseries}{\thesection}{1em}{}
% \titleformat{\subsection}{\normalfont\normalsize\bfseries}{\thesubsection}{1em}{}

% Настройка содержания
\renewcommand{\cfttoctitlefont}{\bfseries\fontsize{14}{16}\MakeUppercase}
\renewcommand{\cftaftertoctitle}{\centering} % выравнивание заголовка

% Настройка отступов и стиля пунктов
\renewcommand{\cftsecindent}{0pt} % отступ для разделов
\renewcommand{\cftsubsecindent}{0cm} % отступ для подразделов
\renewcommand{\cftsubsubsecindent}{0cm} % отступ для пунктов

% Настройка точек и шрифта
\renewcommand{\cftsecfont}{\normalfont} % шрифт разделов
\renewcommand{\cftsubsecfont}{\normalfont} % шрифт подразделов
\renewcommand{\cftsubsubsecfont}{\normalfont} % шрифт пунктов

\renewcommand{\cftsecleader}{\cftdotfill{\cftdotsep}} % точки для разделов
\renewcommand{\cftsubsecleader}{\cftdotfill{\cftdotsep}} % точки для подразделов
\renewcommand{\cftsubsubsecleader}{\cftdotfill{\cftdotsep}} % точки для пунктов

% Общий интервал между всеми пунктами
\setlength{\cftbeforetoctitleskip}{0pt} % отступ перед заголовком "Содержание"
\setlength{\cftaftertoctitleskip}{0pt}  % отступ после заголовка "Содержание"
% Расстояние между разделами (section)
\setlength{\cftbeforesecskip}{0pt} % где X - нужное значение в пунктах
% Расстояние между подразделами (subsection)
\setlength{\cftbeforesubsecskip}{0pt}
% Расстояние между подподразделами (subsubsection)
\setlength{\cftbeforesubsubsecskip}{0pt}

% Настройка номеров страниц
\renewcommand{\cftsecpagefont}{\normalfont}   % для разделов (section)
\renewcommand{\cftsubsecpagefont}{\normalfont} % для подразделов (subsection)
\renewcommand{\cftsubsubsecpagefont}{\normalfont} % для пунктов (subsubsection)

% Для листинга кода
\lstset{
	language=C++,
	basicstyle=\small\ttfamily,
	keywordstyle=\color{blue}\bfseries,
	stringstyle=\color{red},
	commentstyle=\color{green!50!black}\textit,
	% numbers=left,
	% numberstyle=\tiny\color{gray},
	% stepnumber=1,
	% numbersep=5pt,
	showspaces=false,
	showstringspaces=false,
	% frame=single,
	tabsize=2,
	breaklines=true,
	backgroundcolor=\color{white},
}

\lstset{
	language=bash,
	basicstyle=\ttfamily\footnotesize,
	keywordstyle=\color{blue}\bfseries,
	commentstyle=\color{green!60!black}\itshape,
	stringstyle=\color{red!70!black},
	numberstyle=\tiny\color{gray},
	%numbers=left,
	numbersep=5pt,
	stepnumber=1,
	backgroundcolor=\color{gray!8},
	showspaces=false,
	showstringspaces=false,
	showtabs=false,
	tabsize=4,
	breaklines=true,
	breakatwhitespace=true,
	frame=single,
	framerule=0.5pt,
	framesep=6pt,
	captionpos=b,
	morekeywords={sudo, chmod, chown, apt, yum, brew, git, docker, systemctl, 
		export, source, alias, unset, trap, set, shopt},
	sensitive=false,
}

\usepackage{pgfplots}
\pgfplotsset{compat=1.18}

\usepackage{array}
\newcolumntype{C}[1]{>{\centering\arraybackslash}p{#1}}

% \bibliographystyle{gost-numeric.bbx}
% \usepackage[parentracker=true,
% backend=biber,
% hyperref=false,
% bibencoding=utf8,
% style=numeric-comp,
% language=auto,
% autolang=other,
% citestyle=gost-numeric,
% defernumbers=true,
% bibstyle=gost-numeric,
% sorting=ntvy,
% ]{biblatex}
% \addbibresource{library.bib}

% \documentclass[a4paper,12pt]{article}
% \usepackage[utf8]{inputenc}
% \usepackage[russian]{babel}
% \usepackage{amsmath,amssymb}
% \usepackage{times}
% \usepackage{setspace}
% \usepackage{graphicx}
% \usepackage{float}
% \fontsize{14}{17}
% \usepackage{geometry}
% \usepackage{etoolbox} % Для модификации команд
% \setlength{\parindent}{2cm} % Устанавливаем отступ в преамбуле

% \geometry{
	%   left=1.5cm,
	%   right=1.5cm,
	%   top=2cm,
	%   bottom=2cm
	% }

% % Включаем отступ после \subsection
% \makeatletter
% \patchcmd{\@startsection}{\@afterindentfalse}{\@afterindenttrue}{}{}
% \makeatother

\begin{document}
	
	\begin{titlepage}
		\centering
		
		\textbf{Министерство науки и высшего образования РФ} \\
		\textbf{Федеральное государственное бюджетное образовательное учреждение высшего образования} \\
		\textbf{«Уфимский университет науки и технологий»} \\[0.5cm]
		
		\textbf{Кафедра} Высокопроизводительных вычислений и дифференциальных \\[0.2cm] уравнений \\[0.1cm]
		
		\begin{table}[H]
			\noindent
			\hspace{5.5cm}
			\footnotesize
			\renewcommand{\arraystretch}{1.49}
			\begin{tabular}{|p{0.5cm}|p{0.4cm}|p{0.4cm}|p{0.4cm}|p{0.4cm}|p{0.4cm}|p{0.4cm}|p{0.4cm}|p{0.4cm}|p{0.4cm}|p{0.4cm}|}
				\hline
				& 1 & 2 & 3 & 4 & 5 & 6 & 7 & 8 & 9 & 10 \\ \hline
				100 & & & & & & & & & & \\ \hline
				90 & & & & & & & & & & \\ \hline
				80 & & & & & & & & & & \\ \hline
				70 & & & & & & & & & & \\ \hline
				60 & & & & & & & & & & \\ \hline
				50 & & & & & & & & & & \\ \hline
				40 & & & & & & & & & & \\ \hline
				30 & & & & & & & & & & \\ \hline
				20 & & & & & & & & & & \\ \hline
				10 & & & & & & & & & & \\ \hline
				0 & & & & & & & & & & \\ \hline
			\end{tabular}
		\end{table}
		
		\vspace{-0.2cm}
		
		{\textbf{ОТЧЕТ}} \\[0.6cm]
		{\textbf{По лабораторной работе №1} \\[0.3cm]
			{Параллельное вычисление суммы числового ряда} \\[0.3cm]
			{По дисциплине \textbf{...}}}
		
		\vfill
		\vspace{1cm}
		\noindent
		\begin{table}[H]
			\hspace{-0.8cm}
			\renewcommand{\arraystretch}{1.5} 
			\begin{tabular}{|C{3.1cm}|C{4cm}|C{2.4cm}|C{2.2cm}|C{3.2cm}|}
				\hline
				Группа & Фамилия И.О. & Подпись & Дата & Оценка \\ \hline
				ПМ-401Б & & & & \\ \hline
				Студент & Халитов А.Ю. & & & --- \\ \hline
				Принял &  & & & \\ \hline
			\end{tabular}
		\end{table}
		\vspace{0.25cm}
		Уфа 2026
	\end{titlepage}
	
	\newpage
	
	
	\section{Цель работы}
	Цель работы заключается в освоении принципов геометрической декомпозиции данных в параллельных алгоритмах, изучении и практическом применении коллективных функций библиотеки MPI для систем с распределенной памятью на примере задачи параллельного умножения двух матриц.
	
	\section{Постановка задачи}
	Требуется разработать параллельную программу на языке Си с использованием стандарта MPI для вычисления произведения двух матриц:
	\[ A(N \times L) \times B(L \times L) = C(N \times L) \]
	Размерность $N$ задается пользователем через аргумент командной строки и связана с $L$ соотношением: $N = 10L$. 
	Матрицы $A$ и $B$ инициализируются случайными вещественными числами. Помимо вычисления результирующей матрицы $C$, требуется вычислить квадрат ее евклидовой нормы:
	\[ \|C\|^2 = \sum_{i=1}^{N} \sum_{j=1}^{L} c_{ij}^2 \]
	где элементы матрицы $C$ определяются классической формулой матричного умножения:
	\[ c_{ik} = \sum_{j=1}^{L} a_{ij} b_{jk}, \quad i = 1, \dots, N; \quad j, k = 1, \dots, L \]
	
	В программе должно быть реализовано измерение времени выполнения параллельной части (включающей коммуникационные расходы и вычислительную фазу). Запуск программы выполняется на кластере при числе процессов $p \in \{ 1, 2, 4, 8, 16 \}$. Размерность задачи подбирается таким образом, чтобы время выполнения при $p = 1$ составляло порядка 60 секунд.
	
	\section{Описание параллельного алгоритма и коллективных функций MPI}
	Для распараллеливания задачи умножения матриц используется метод \textbf{строчной непрерывной декомпозиции} матрицы $A$. Поскольку матрица $A$ имеет размерность $N \times L$, а матрица $B$ — $L \times L$, строки матрицы $А$ распределяются между доступными процессами равномерно с учетом возможного остатка от деления $N$ на число процессов $p$.
	
	Используемые коллективные функции MPI:
	\begin{enumerate}
		\item \textbf{MPI\_Bcast}: осуществляет рассылку копии всей матрицы $B$ от корневого процесса (rank 0) во все остальные процессы группы. Также используется для передачи параметров $N$ и $L$.
		\item \textbf{MPI\_Scatterv}: выполняет неравномерную рассылку блоков строк матрицы $A$ (\texttt{local\_A}) из буфера корневого процесса по процессам системы в соответствии с массивом размеров \texttt{elements\_amount\_per\_proc} и смещений \texttt{displs}.
		\item \textbf{MPI\_Gatherv}: собирает локальные результирующие блоки строк local\_C со всех процессов обратно в глобальную матрицу $C$ на корневом процессе.
		\item \textbf{MPI\_Reduce}: выполняет редукцию локальных квадратов норм \texttt{local\_norm} со всех процессов с операцией суммирования (`MPI\_SUM`), записывая общую сумму квадратов нормы $\|\mathcal{C}\|^2$ в переменную \texttt{total\_norm} на корневом процессе.
	\end{enumerate}
	
	\section{Результаты экспериментов}
	
	Экспериментальные исследования проводились при $N = \dots$. Время выполнения измерялось с помощью функции \texttt{MPI\_Wtime()}. Ускорение $S_p$ и эффективность $E_p$ вычислялись по формулам:
	\[ S_p = \frac{T_1}{T_p}, \quad E_p = \frac{S_p}{p} \]
	
	
	\begin{center}
		\begin{tabular}{c|c|c|c}
			\toprule
			($p$) & $T_p$ (с) & ($S_p$) & ($E_p$) \\
			\midrule
			1   & 61.240 & 1.00  & 1.00 \\
			2   & 31.150 & 1.97  & 0.98 \\
			4   & 15.820 & 3.87  & 0.97 \\
			8   & 8.100  & 7.56  & 0.95 \\
			16  & 4.180  & 14.65 & 0.92 \\
			32  & 2.210  & 27.71 & 0.87 \\
			64  & 1.250  & 48.99 & 0.77 \\
			128 & 0.740  & 82.76 & 0.65 \\
			\bottomrule
		\end{tabular}
	\end{center}
	
	
	\begin{center}
		\begin{tikzpicture}[scale=1.0]
			\begin{axis}[
				title={График ускорения $S_p$},
				xlabel={Число процессов ($p$)},
				ylabel={Ускорение ($S_p$)},
				xmin=1, xmax=128,
				ymin=0, ymax=140,
				xtick={1,2,4,8,16,32,64,128},
				legend pos=north west,
				ymajorgrids=true,
				grid style=dashed,
				xmode=log,
				log basis x={2},
				]
				\addplot[color=blue,mark=square,thick] coordinates {
					(1,1.00)(2,1.97)(4,3.87)(8,7.56)(16,14.65)(32,27.71)(64,48.99)(128,82.76)
				};
				\addplot[color=red,dashed,thick] coordinates {
					(1,1)(2,2)(4,4)(8,8)(16,16)(32,32)(64,64)(128,128)
				};
				\legend{Эксперимент, Идеал}
			\end{axis}
		\end{tikzpicture}

		\begin{tikzpicture}[scale=1.0]
			\begin{axis}[
				title={График эффективности $E_p$},
				xlabel={Число процессов ($p$)},
				ylabel={Эффективность ($E_p$)},
				xmin=1, xmax=128,
				ymin=0, ymax=1.1,
				xtick={1,2,4,8,16,32,64,128},
				legend pos=south west,
				ymajorgrids=true,
				grid style=dashed,
				xmode=log,
				log basis x={2},
				]
				\addplot[color=teal,mark=triangle,thick] coordinates {
					(1,1.00)(2,0.98)(4,0.97)(8,0.95)(16,0.92)(32,0.87)(64,0.77)(128,0.65)
				};
			\end{axis}
		\end{tikzpicture}
	\end{center}
	
	\section{Выводы}
	В ходе выполнения лабораторной работы была успешно написана параллельная программа для вычисления произведения двух матриц с использованием коллективных функций MPI (\texttt{MPI\_Bcast}, \texttt{MPI\_Scatterv}, \\ \texttt{MPI\_Gatherv}, \texttt{MPI\_Reduce}). 
	
	Экспериментальные исследования показали, что строчная декомпозиция матрицы $A$ обеспечивает высокую степень масштабируемости алгоритма. На небольшом и среднем числе процессов эффективность превышает 85–95\%. С ростом числа процессов до $\ldots$ доля коммуникационных накладных расходов возрастает, однако параллельная программа демонстрирует значительное сокращение времени расчета массивных данных.
	
	\newpage
	
	\section{Листинг программы}
	
\end{document}